\section{Proof of Theorem~\ref{thm:exact}}\label{sec:proof}

Let $D_*$ be as in Theorem~\ref{thm:half} and $K_0$ as in Section~\ref{sec:poles}.
By~\eqref{eq:H},
\begin{equation}\label{eq:Hlower}
  H(r)\;\ge\;2^{r-1}+(r-2)\,2^{r/2-1}\qquad\text{for all }r\ge3 .
\end{equation}
Choose $R\ge8$ so large that the following holds for every $r\ge R$:
Proposition~\ref{prop:threecut}, Theorem~\ref{thm:fourcut}, Theorem~\ref{thm:girth} with
$g_0=D_*-1$, and Theorem~\ref{thm:half} apply to the rank $r$; and $r-2>10K_0^{2}$ and
$2\log(2r)+2\le3\log r$.

Let $r\ge R$. A ring $R_r$ has $H(r)$ cycles by Corollary~\ref{cor:ringmax}, so $\psi(r)\ge H(r)$.
By Lemma~\ref{lem:cubic}, some connected cubic multigraph $G$ of rank $r$ has $\psi(r)$ cycles. We
show that $G$ is a ring or that $c(G)\le2^{r-1}+5K_0^{2}\,2^{r/2}$. In the first case
$c(G)\le H(r)$ by Corollary~\ref{cor:ringmax}. In the second case $c(G)<H(r)$ by~\eqref{eq:Hlower},
because $r-2>10K_0^{2}$. So $\psi(r)\le H(r)$.

\step{Case 1: $G$ has a loop $e$.} Every cycle of $G$ other than $e$ is a cycle of $G-e$, which has
rank $r-1$. So $c(G)\le1+(2^{r-1}-1)$.

\step{Case 2: $G$ has no loop, and it has a cyclic cut of size at most three.} Then
$c(G)\le2^{r-1}+5K_0^{2}\,2^{r/2}$ by Proposition~\ref{prop:threecut}.

\step{Case 3: $G$ has no loop, it has a cyclic cut of size four, and it has no smaller cyclic
cut.} Then $G$ is a ring or $c(G)<2^{r-1}$, by Theorem~\ref{thm:fourcut}.

In the other cases $G$ has no loop and no cyclic cut of size at most four. Let $g$ be its girth.

\step{Case 4: $g<D_*$.} Then $c(G)<2^{r-1}$ by Theorem~\ref{thm:girth}.

\step{Case 5: $g\ge D_*$.} Let $W$ be a shortest cycle, and put $P=V(W)$. As in the proof of
Corollary~\ref{cor:cycles}, $|P|=g\le2\log(2r)+2\le3\log r$. The graph $G[P]$ is connected, and
$G-P$ has girth at least $D_*$. So $q(G,P)<\frac25$ by Theorem~\ref{thm:half}.

Let $Z$ be a cycle of $G$ that meets $P$. Its edges in $G-P$ form a proper subset of $Z$, hence a
forest. So at most $q(G,P)\,2^{r}<\frac25\,2^{r}$ cycles meet $P$. The other cycles are cycles of
$G-P$, and there are at most $2^{\beta(G-P)}$ of them. The cycle $W$ has no chord. So $G-P$ arises
from $G$ by the deletion of $g$ vertices and $2g$ edges, and
\[
  \beta(G-P)=r-1-g+\comp(G-P).
\]
The graph $G$ has no cut of size one, because such a cut is cyclic (Section~\ref{sec:cuts}). So
every component of $G-P$ contains the ends of at least two of the $g$ edges that leave $P$. Hence
$\comp(G-P)\le g/2$ and $\beta(G-P)\le r-1-g/2\le r-4$, because $g\ge7$. Therefore
$c(G)<\bigl(\frac25+\frac1{16}\bigr)2^{r}<2^{r-1}$. \qed
